\subsection*{Проверка переноса по температуре}
При \(q_c=0{,}37\) ноль \(\max_{\Delta\mu}\ln R\) найден совместным перевзвешиванием по энергии и составу. После снятия зонтичного потенциала при \((\beta_0,\Delta\mu_0)\) логарифм поправки к весу кадра равен \(-(\beta-\beta_0)E+(\beta\Delta\mu-\beta_0\Delta\mu_0)k\). В совместном расчёте MBAR нормировки всех зонтиков разных температур оцениваются по исходным кадрам. Блочные интервалы условны по размеру, траекториям и разделителю.
@THERMALTABLE@
На 2304 атомах первая и вторая половины траекторий дают 1242,7 и 1244,6~К. Изменение \(q_c\) с 0,32 до 0,47 сдвигает \(T_*\) на 0,17~К; на 288 атомах сдвиг около 7~К. На 768 атомах независимые переносы с 1225, 1250 и 1273,15~К дают 1246,6, 1243,1 и 1247,0~К.

На 2304 атомах локальный температурный наклон \(\max_{\Delta\mu}\ln R\) при 1245~К и \(q_c=0{,}47\) даёт 1243,52~К [1241,83; 1245,18]. Он проверяет нелинейное перевзвешивание на тех же траекториях другим способом. В совместном MBAR учёт каждого четвёртого вместо каждого восьмого кадра меняет \(T_*\) на 0,25--0,28~К.
